\section{Hiệu năng mạng trong public cloud}

CloudCmp đề xuất benchmark nhiều chiều để so sánh public cloud và chỉ ra rằng
kết luận phụ thuộc workload, metric và thời điểm đo \cite{cloudcmp}. Cách tiếp
cận này hỗ trợ lựa chọn phương pháp của đồ án: không dùng một metric duy nhất
để kết luận toàn bộ giải pháp kết nối.

Wang và Ng đo ảnh hưởng của virtualization lên network performance của Amazon
EC2, cho thấy biến thiên và overhead của nền tảng ảo hóa cần được xem xét khi
giải thích latency và throughput \cite{wang-ec2-network}. Nghiên cứu đó tạo
bối cảnh cho việc lưu CPU, retransmission và nhiều lượt đo, nhưng không trực
tiếp so sánh VPC Peering với TGW trên cùng endpoint.

\section{Tài liệu AWS về kết nối liên VPC}

AWS mô tả VPC Peering là kết nối point-to-point, không bắc cầu; TGW là router
khu vực kết nối VPC và mạng ngoài AWS qua attachment
\cite{aws-vpc-peering,aws-tgw-overview}. Tài liệu route, quota và pricing xác
định đặc tính quản trị, scale và cost model, nhưng không công bố một giá trị
RTT hoặc TCP goodput áp dụng cho mọi workload.

\section{Tài liệu về EC2, ENA và công cụ benchmark}

Tài liệu EC2 network bandwidth giải thích giới hạn của instance
\cite{aws-ec2-bandwidth}. Tài liệu ENA mô tả ý nghĩa của allowance counter
\cite{aws-ena-metrics}. iperf3 cung cấp JSON gồm byte, thời gian, bitrate phía
nhận và retransmission; công cụ đo end-to-end nên không cô lập riêng một nút
mạng \cite{iperf3-docs}.

\section{Vị trí của đồ án}

Đồ án sử dụng cùng bốn EC2, cùng Availability Zone, private IP, MTU và tham số
tải; chỉ chuyển 12 route giữa Peering và TGW. Ngoài RTT và TCP single-flow,
nghiên cứu còn quan sát fan-in, CPU và ENA, đồng thời giữ raw artifact để tính
lại. Thiết kế này giảm một số confounder so với hai environment độc lập nhưng
không loại bỏ biến thiên host, không thay thế benchmark đa region và không tạo
cơ sở để tuyên bố đây là nghiên cứu đầu tiên về chủ đề này.
